\chapter{Balls, Bins, and Random Graphs}

\paragraph{Exercise 5.2} Let us start by calculating how many people we would
need to have in a room before it was more likely than not that two had the
same last four digits. It is easier to think about the configuration where
people do not share a password than about configurations where some two people
do. The first person in the room has a password. The probability that the second
person has a different password is $\left(1- \frac{1}{4^{10}}\right)$. The
probability that the third person in the room then has a password dofferent from
the first two, given that the first two people have different passwords, is
$\left(1- \frac{2}{4^{10}}\right)$. Continuing on, the probability that the
$kth$ person in the room has a different password than the first $k-1$,
assuming that the first $k-1$ have different passwords, is
$\left(1- \frac{k-1}{4^{10}}\right)$. So the probability that $n$ people all
have different passwords is
\[
  \prod_{j=1}^{n-1} \left( 1 - \frac{j}{4^{10}}\right).
\]
Using that $1 - \frac{k}{n} \approx e^{-\frac{k}{n}}$ when $k \ll n$, we see that
if $n$ is small compared to $4^{10}$ then
\[
  \prod_{j=1}^{n-1} \left( 1 - \frac{j}{4^{10}}\right)
  /approx \prod_{j=1}^{n-1} e^{-\frac{j}{4^{10}}}
  = \exp \left\{ - \sum_{j=1}^{n-1} \frac{j}{4^{10}} \right\}
  = e^{-\frac{(n-1)n}{2 \cdot 4^{10}}}
  \approx e^{-\frac{n^2}{2 \cdot 4^{10}}}.
\]
Hence the value for $n$ at which the probability that $n$ people all have
different passwords is smaller than $1/2$ is approxuimately given by the
unequality
\[
  e^{-\frac{n^2}{2 \cdot 4^{10}}} < \frac{1}{2}
  \Leftrightarrow e^{\frac{n^2}{2 \cdot 4^{10}}} > 2
  \Leftrightarrow \frac{n^2}{2 \cdot 4^{10}} > \ln(2)
  \Leftrightarrow n > \sqrt{2 \cdot 4^{10} \cdot \ln(2)}.
\]
Therefore, with 1206 people in a room the probability that two of them have
the same password is greater than $1/2$. \\
Similarly one can argue that $\sqrt{2 \cdot 9^{10} \cdot \ln(2)} \approx 69,524$
numbers can be issued before it will be more likely than not that there is a
duplicate number. \\
If Social Security numbers had 13 digits we would still nee 1206 people in a room
before it is more likely than not that two have the same password. However, we
could issue $\sqrt{2 \cdot 13^{10} \cdot \ln(2)} \approx 437,164$ numbers before
it is more likely than not that there is a duplicate number.

\paragraph{Exercise 5.6}
Let $i \in \mathbb{N}$ and let $X_i$ be the event that there are $i$ errors on a
page, that is $X = i$. Assume that $X_n$ occurs. Then $Y$ is a binomial random
variable with parameters $n$ and $p$, $Y \sim B(n,p)$. Similarly, $Z$ is a binomial
random variable with parameters $n$ and $(1-p)$, $Z \sim B(n,1-p)$.
Thus, for any $y,z \in \mathbb{N}$,
\[ \pr\left(Y = y | X = n\right) = \binom{n}{y} p^y (1-p)^{n-y}, \]
and
\[ \pr\left(Z = z | X = n\right) = \binom{n}{z} (1-p)^z p^{n-z}. \]
Since $X_0, X_1, X_2, ...$ are mutually disjoint events, such that
$\bigcup_{i \in \mathbb{N}} X_i = \Omega$, we can apply the law of total probability
to compute the probabilities $\pr(Y = y)$ and $\pr(Z = z)$:
\begin{align*}
  \pr(Y = y)
    &= \sum_{n = y}^{\infty} \pr(Y = y | X = n) \cdot \pr(X = n) \\
    &= \sum_{n = y}^{\infty} \binom{n}{y} p^y (1-p)^{n-y} \cdot \frac{e^{-\mu}\mu^n}{n!} \\
    &= e^{-\mu} p^y \sum_{n = y}^{\infty} \frac{n!}{y!(n-y)!} \cdot (1-p)^{n-y} \cdot \frac{\mu^n}{n!} \\
    &= \frac{e^{-\mu} p^y}{y!} \sum_{n = y}^{\infty} \frac{(1-p)^{n-y}}{(n-y)!} \cdot \mu^n \\
    &= \frac{e^{-\mu} p^y}{y!} \sum_{n = y}^{\infty} \frac{(1-p)^{n-y}}{(n-y)!} \cdot \mu^n \cdot \frac{\mu^y}{\mu^y} \\
    &= \frac{e^{-\mu} p^y \mu^y}{y!} \sum_{n = y}^{\infty} \frac{(1-p)^{n-y}}{(n-y)!} \cdot \mu^{n-y} \\
    &= \frac{e^{-\mu} p^y \mu^y}{y!} \sum_{n = y}^{\infty} \frac{\left((1-p)\mu\right)^{n-y}}{(n-y)!} \\
    &= \frac{e^{-\mu} p^y \mu^y}{y!} \sum_{n = 0}^{\infty} \frac{\left((1-p)\mu\right)^{n}}{n!} \\
    &= \frac{e^{-\mu} p^y \mu^y}{y!} \cdot e^{(1-p)\mu} \\
    &= \frac{e^{-\mu + \mu - \mu p} p^y \mu^y}{y!} \\
    &= \frac{e^{- \mu p} (\mu p)^y}{y!}, \\
\end{align*}
where we have used the Taylor expansion $e^x = \sum_{j=0}^\infty \frac{x^j}{j!}$.
Therefore, $Y$ is a Poisson random variable with parameter $\mu p$. \\
Similarly, one has
\begin{align*}
  \pr(Z = z)
    &= \sum_{n = z}^{\infty} \pr(Z = z | X = n) \cdot \pr(X = n) \\
    &= \sum_{n = z}^{\infty} \binom{n}{z} (1-p)^z p^{n-z} \cdot \frac{e^{-\mu}\mu^n}{n!} \\
    &= e^{-\mu} (1-p)^z \sum_{n = z}^{\infty} \frac{n!}{z!(n-z)!} \cdot p^{n-z} \cdot \frac{\mu^n}{n!} \\
    &= \frac{e^{-\mu} (1-p)^z}{z!} \sum_{n = y}^{\infty} \frac{p^{n-z}}{(n-z)!} \cdot \mu^n \cdot \frac{\mu^z}{\mu^z} \\
    &= \frac{e^{-\mu} \left(\mu(1-p)\right)^z}{z!} \sum_{n = z}^{\infty} \frac{(p\mu)^{n-y}}{(n-y)!} \\
    &= \frac{e^{-\mu} \left(\mu(1-p)\right)^z}{z!} \sum_{n = 0}^{\infty} \frac{\left(p\mu\right)^{n}}{n!} \\
    &= \frac{e^{-\mu} \left(\mu(1-p)\right)^z}{z!} \cdot e^{p\mu} \\
    &= \frac{e^{-\mu + p\mu} \left(\mu(1-p)\right)^z}{z!} \\
    &= \frac{e^{-\mu(1-p)} \left(\mu(1-p)\right)^z}{z!}. \\
\end{align*}
Consequently, $Z$ is a Poisson random variable with parameter $\mu (1-p)$. \\
Additionally, we want to show that $Y$ and $Z$ are independent. Let $y,z \in \mathbb{N}$.
Then,
\begin{align*}
  \pr(Y = y \cap Z = z)
    &= \pr(Y = y \cap X = y + z) \\
    &= \pr(Y = y | X = y + z) \cdot \pr(X = y + z) \\
    &= \binom{y+z}{y} p^y (1-p)^{(y+z)-y} \cdot \frac{e^{-\mu}\mu^{y+z}}{(y+z)!} \\
    &= \frac{(y+z)!}{y!(y+z-y)!} p^y (1-p)^z \cdot \frac{e^{-\mu(p + (1-p))}\mu^y \cdot \mu^z}{(y+z)!} \\
    &= \frac{e^{-\mu p}(\mu p)^y}{y!} \cdot \frac{e^{-\mu(1-p)} \left(\mu (1-p)\right)^z}{z!} \\
    &= \pr(Y = y) \cdot \pr(Z = z).
\end{align*}
Hence, $Y$ and $Z$ are independent.
